\section{The semiregularity target}\label{sec:semireg}

Markman's construction at $n=2$ deforms a sheaf across a family and controls
the obstruction by a semiregularity argument in the sense of Bloch and of
Buchweitz and Flenner \cite{Blo72,BF03}. As for the secant construction of
\Cref{sec:numerology}, one may ask whether a dimension count blocks this route
for $n\ge4$. We compute the target of the semiregularity map and show that it
grows too fast for the count to impose any such constraint, so that neither
route meets a numerical obstruction. The constraint that does exist is
structural. In \Cref{ssec:rankzero} we show that the annihilator of the Weil
line in $H^{0,2}(A)$ has dimension $n^{2}$, and that it prevents every object
whose Chern character is concentrated on the Weil line from being
semiregular as soon as $H^{2}(\cO_{A})$ acts on the object faithfully. In
\Cref{ssec:tangent} we identify this annihilator with the tangent space to the
Weil family.

\subsection{The dimension of the target}

\begin{lemma}[Hodge numbers of an abelian variety]\label{lem:hodgenumbers}
For an abelian variety $A$ of dimension $g$ one has
$h^{p,q}(A)=\binom{g}{p}\binom{g}{q}$.
\end{lemma}

\begin{proof}
By \eqref{eq:wedge} and the compatibility of the Hodge decomposition with the
exterior algebra structure,
$H^{p,q}(A)=\bigwedge^{p}H^{1,0}\otimes\bigwedge^{q}H^{0,1}$, and
$\dim H^{1,0}=\dim H^{0,1}=g$.
\end{proof}

\begin{proposition}[The semiregularity target]\label{prop:target}
Let $A$ be an abelian variety of dimension $g=2n$. The target of the
semiregularity map, namely $\bigoplus_{q}H^{q+2}(A,\Omega^{q}_{A})$, has
dimension
\[
  \sum_{q\ge0}h^{q,q+2}(A)=\binom{4n}{2n-2}.
\]
\end{proposition}

\begin{proof}
By \Cref{lem:hodgenumbers} with $g=2n$,
\[
  \sum_{q\ge0}h^{q,q+2}
  =\sum_{q\ge0}\binom{2n}{q}\binom{2n}{q+2}
  =\sum_{q\ge0}\binom{2n}{q}\binom{2n}{2n-q-2}
  =\binom{4n}{2n-2},
\]
the last step being the Vandermonde convolution.
\end{proof}

\begin{corollary}[The dimension count]\label{cor:countnotbind}
The target dimension $\binom{4n}{2n-2}$ takes the values
$28,495,8008,125970,1961256,30421755$ for $n=2,\dots,7$, growing by a factor
exceeding ten at each step. Any source whose dimension grows polynomially in
$n$, for instance the quadratic quantity $\binom{2n}{2}=n(2n-1)$, which takes
the values $6,15,28,45,66,91$ over the same range, is eventually negligible
against it. In particular, for large $n$ a comparison of the two dimensions
cannot obstruct the existence of a semiregular sheaf; the inequality it
produces points the other way.
\end{corollary}

\begin{remark}[Consistency at $n=2$]\label{rem:markmanconsistent}
At $n=2$ the target has dimension $28$, and the sheaf produced by Markman has
$\dim\Ext^{2}=12$, well below it, which is consistent with the construction
succeeding there. Nothing in the numbers changes character between $n=2$ and
$n=4$, so if the construction cannot be carried out for larger $n$, the reason
is not this comparison.
\end{remark}

\begin{remark}[The remaining requirement]\label{rem:semiregneed}
\Cref{cor:countnotbind} does not show that the semiregularity route works.
It removes one apparent obstruction and leaves the actual requirement in
place, namely the construction of a sheaf on $A$ whose Chern character has a
nonzero component along the Weil line. By \Cref{thm:character} such a sheaf
cannot be $K$-equivariant with the norm character, so it cannot be produced
by any functorial construction from $K$-invariant data on $\Av$; Markman's
sheaf at $n=2$ is not produced in this way. In this sense the problem for
$n\ge3$ is one of breaking equivariance, and the dimensions play no role in
it.
\end{remark}

\subsection{The annihilator of the Weil line and objects of rank zero}
\label{ssec:rankzero}

By \Cref{rem:semiregneed}, what remains is to construct a sheaf whose Chern
character has a component along the Weil line. The most economical way to
arrange this is to make the Chern character lie \emph{entirely} on the Weil
line, that is, to take an object of rank zero with every $\ch_{k}$ vanishing
except $\ch_{n}$, as the construction of \Cref{sec:construction} does. We show
that this choice is incompatible with semiregularity whenever $H^{2}(\cO_{A})$
acts faithfully on the object, for a reason that depends only on the position
of the Chern character and not on how the object is built.

We first recall the semiregularity map, on which the rest of this subsection
turns. For a coherent sheaf $\cE$ on a smooth projective variety $A$ the semiregularity
map is defined through the Atiyah class
$\mathrm{At}(\cE)\in\Ext^{1}(\cE,\cE\otimes\Omega^{1}_{A})$: a class
$\xi\in\Ext^{2}(\cE,\cE)$ is sent to the trace of
$\xi\cdot\exp(\mathrm{At}(\cE))$, and the component of that trace in
$H^{q+2}(A,\Omega^{q}_{A})=H^{q,q+2}(A)$ is the $q$-th coordinate of
$\sigma_{\cE}(\xi)$. Bloch introduced the map for local complete
intersections \cite{Blo72}, and Buchweitz and Flenner extended it to
arbitrary coherent sheaves and perfect complexes \cite{BF03}. If $A$ moves in
a family $\cA\to S$, the obstruction to extending $\cE$ to first order along
$v\in T_{s}S$ is a class $\mathrm{ob}_{v}(\cE)\in\Ext^{2}(\cE,\cE)$, and the
theorem of Buchweitz and Flenner identifies $\sigma_{\cE}(\mathrm{ob}_{v}(\cE))$
with the derivative of $\ch(\cE)$ in the direction $v$, projected off the
diagonal of the Hodge decomposition. Hence if $\ch(\cE)$ stays of Hodge type
along $S$, then $\sigma_{\cE}(\mathrm{ob}_{v}(\cE))=0$, and if $\sigma_{\cE}$ is
also injective, then $\mathrm{ob}_{v}(\cE)=0$ and $\cE$ deforms.
Injectivity thus converts a statement about the Chern character into a
statement about the sheaf, and it is the one hypothesis in the argument that
sees more than $\ch(\cE)$.

Two special cases are all that is needed below. For a line bundle $L$ one has
$\Ext^{2}(L,L)=H^{2}(\cO_{A})=H^{0,2}(A)$, and $\sigma_{L}$ is multiplication
by $e^{c_{1}(L)}$. For a direct sum, $\Ext^{2}$ is the sum of the
$\Ext^{2}(\cE_{i},\cE_{j})$, and $\sigma$ acts on each diagonal block
$H^{2}(\cO_{A})\cdot\id_{\cE_{i}}$ by multiplication by the Chern character of
that summand. Both formulae follow at once from the definition by the trace.

\begin{lemma}[The annihilator of the Weil line in $H^{0,2}$]
\label{lem:annihilator}
Let $A$ be of $(K,-1,n)$-Weil type and let $\omega\in\HW(A)\otimes\RR$ be
any nonzero real element of the Weil line. Then the multiplication map
\[
  H^{0,2}(A)\longrightarrow H^{n,n+2}(A),\qquad z\longmapsto z\cdot\omega,
\]
has kernel of dimension exactly $n^{2}$, spanned by the products $u\wedge v$
with $u\in V_{+}^{0,1}$ and $v\in V_{-}^{0,1}$. In particular the kernel is
the same subspace for every $\omega$ on the line.
\end{lemma}

\begin{proof}
Write $\omega=a\,\alpha_{+}+b\,\alpha_{-}$ with $\alpha_{\pm}$ generating
$\bigwedge^{2n}V_{\pm}$, which is possible because
$\HW(A)\otimes\CC=\bigwedge^{2n}V_{+}\oplus\bigwedge^{2n}V_{-}$ (see the proof
of \Cref{prop:weilline}). Complex conjugation interchanges $V_{+}$ and
$V_{-}$, so $\bbar{\alpha}_{+}=\mu\,\alpha_{-}$ for some $\mu\ne0$, and
comparing $\omega$ with $\bbar{\omega}=\omega$ gives $b=\bbar{a}\,\mu$. Since
$\omega\ne0$, the coefficients $a$ and $b$ are therefore both nonzero. As
$H^{0,1}=V_{+}^{0,1}\oplus V_{-}^{0,1}$ with each summand of dimension $n$,
the space $H^{0,2}=\bigwedge^{2}H^{0,1}$ decomposes as
\[
  H^{0,2}
  =\textstyle\bigwedge^{2}V_{+}^{0,1}
  \;\oplus\;\bigl(V_{+}^{0,1}\otimes V_{-}^{0,1}\bigr)
  \;\oplus\;\bigwedge^{2}V_{-}^{0,1},
\]
of dimensions $\binom{n}{2}$, $n^{2}$ and $\binom{n}{2}$.
For $u\in V_{+}^{0,1}$ and $v\in V_{-}^{0,1}$ the mixed product $u\wedge v$
satisfies $u\wedge v\wedge\alpha_{+}=0$, because $u\wedge\alpha_{+}$ lies in
$\bigwedge^{2n+1}V_{+}=0$, and $u\wedge v\wedge\alpha_{-}=0$ for the same
reason with $v$. So the middle summand lies in the kernel.

Conversely, let $z=z_{++}+z_{+-}+z_{--}$ be in the kernel. For the same
reason $z_{++}\wedge\alpha_{+}$, $z_{--}\wedge\alpha_{-}$ and
$z_{+-}\wedge\omega$ vanish, so
$z\cdot\omega=a\,z_{--}\wedge\alpha_{+}+b\,z_{++}\wedge\alpha_{-}$. If $n=1$
the outer summands are zero and there is nothing to prove. If $n\ge2$ the two
terms lie in the distinct summands
$\bigwedge^{2n}V_{+}\otimes\bigwedge^{2}V_{-}$ and
$\bigwedge^{2}V_{+}\otimes\bigwedge^{2n}V_{-}$ of
$\bigwedge^{2n+2}V_{\CC}=\bigwedge^{2n+2}(V_{+}\oplus V_{-})$, so each of them
vanishes. Wedging with a generator of $\bigwedge^{2n}V_{+}$ is injective on
$\bigwedge^{2}V_{-}$, and wedging with a generator of $\bigwedge^{2n}V_{-}$ is
injective on $\bigwedge^{2}V_{+}$. Since $a\neq0$ we get $z_{--}=0$, and since
$b\neq0$ we get $z_{++}=0$, so the kernel is the middle summand.
\end{proof}

The reality of $\omega$ is needed: for the non-real element
$\omega=\alpha_{+}$ of $\HW(A)\otimes\CC$ the same computation gives the
kernel $V_{+}^{0,1}\otimes V_{-}^{0,1}\oplus\bigwedge^{2}V_{+}^{0,1}$, of
dimension $n^{2}+\binom{n}{2}$. \Cref{fig:fibres} draws the mixed block that forms the kernel.

\begin{theorem}[The rank-zero obstruction]
\label{thm:rankzero}
Let $A$ be of $(K,-1,n)$-Weil type. Let $\cE_{1}$ and $\cE_{2}$ be coherent
sheaves on $A$ such that
\[
  \ch(\cE_{1})-\ch(\cE_{2})=c\,\omega,\qquad c\neq0,\quad \omega\in\HW(A),
\]
such that $\Ext^{2}(\cE_{1},\cE_{2})=\Ext^{2}(\cE_{2},\cE_{1})=0$, and such
that the map $H^{2}(\cO_{A})\to\Ext^{2}(\cE_{1},\cE_{1})$,
$z\mapsto z\cdot\id_{\cE_{1}}$, is injective, as it is whenever
$\rk\cE_{1}\neq0$. Put $\cE=\cE_{1}\oplus\cE_{2}$. Then the semiregularity
map
\[
  \sigma_{\cE}\colon\ \Ext^{2}(\cE,\cE)\longrightarrow\bigoplus_{q\ge0}H^{q,q+2}(A)
\]
has kernel of dimension at least $n^{2}$. In particular $\cE$ is not
semiregular, for any $n\ge1$.
\end{theorem}

\begin{proof}
The vanishing of the off-diagonal $\Ext^{2}$ groups gives
$\Ext^{2}(\cE,\cE)=\Ext^{2}(\cE_{1},\cE_{1})\oplus\Ext^{2}(\cE_{2},\cE_{2})$,
and each summand receives $H^{2}(\cO_{A})$ through the $\cO_{A}$-module
structure, $z\mapsto z\cdot\id_{\cE_{i}}$. For $z\in H^{2}(\cO_{A})=H^{0,2}(A)$
consider the element $\xi_{z}=z\cdot\id_{\cE_{1}}-z\cdot\id_{\cE_{2}}$. Its
component in $\Ext^{2}(\cE_{1},\cE_{1})$ is $z\cdot\id_{\cE_{1}}$, so by
hypothesis $z\mapsto\xi_{z}$ is injective. The semiregularity map is
compatible with direct sums, and on $H^{0,2}\cdot\id_{\cE_{i}}$ it is
multiplication by the Chern character of $\cE_{i}$, so
\[
  \sigma_{\cE}(\xi_{z})=z\cdot\ch(\cE_{1})-z\cdot\ch(\cE_{2})
  =c\,z\cdot\omega .
\]
If $\omega=0$ this vanishes for every $z$, and
$\dim H^{0,2}=n(2n-1)\ge n^{2}$. Otherwise $\omega$, being rational, is a
nonzero real element of the Weil line, and by \Cref{lem:annihilator} the
product vanishes for every $z$ in an $n^{2}$-dimensional subspace of
$H^{0,2}$. The corresponding $\xi_{z}$ are linearly independent, and hence
$\dim\ker\sigma_{\cE}\ge n^{2}$.

It remains to justify the claim in the statement that the hypothesis on
$\cE_{1}$ holds in positive rank. For a coherent sheaf $\cF$ the map
$H^{2}(\cO_{A})\to\Ext^{2}(\cF,\cF)$, $z\mapsto z\cdot\id_{\cF}$, is the edge
map of the local-to-global spectral sequence, and
$\operatorname{tr}(z\cdot\id_{\cF})=\rk(\cF)\,z$. When $\rk\cF\ne0$ the trace
therefore splits the map up to the factor $\rk\cF$, and the map is
injective; this applies in particular to a direct sum of line bundles. For a
sheaf of rank zero the map can vanish: for a skyscraper sheaf it factors
through $H^{2}$ of the structure sheaf of a point, which is zero. The direct
sum $\cE_{1}\oplus\cE_{2}$ is used only to produce an element of $\Ext^{2}$
whose image under $\sigma$ is $z\cdot(\ch\cE_{1}-\ch\cE_{2})$ rather than
$z\cdot(\ch\cE_{1}+\ch\cE_{2})$.
\end{proof}

\begin{corollary}[The construction of \Cref{sec:construction}]
\label{cor:pteobstructed}
Let $L_{1},\dots,L_{m}$ be line bundles on $A$ with
$\sum_{i}c_{1}(L_{i})=0$ such that the virtual class
$\sum_{i}\bigl([L_{i}]-[L_{i}^{-1}]\bigr)$ has Chern character on the Weil
line and every difference
$\varepsilon\,c_{1}(L_{i})-\varepsilon'c_{1}(L_{j})$ with
$\varepsilon,\varepsilon'\in\{\pm1\}$ and $(\varepsilon,i)\ne(\varepsilon',j)$
is nondegenerate of index different from $2$. Then the sheaf
$\bigoplus_{i}(L_{i}\oplus L_{i}^{-1})$ has semiregularity kernel of
dimension at least $n^{2}$. At $n=3$ the explicit object of
\Cref{sec:construction} has kernel of dimension exactly
$15=n^{2}+n(n-1)$, of which the $n^{2}=9$ dimensions on the odd side are
the ones supplied by \Cref{thm:rankzero}; the remaining $n(n-1)=6$ come from
the even side and are specific to the line-bundle structure.
\end{corollary}

\begin{proof}
Take $\cE_{1}=\bigoplus L_{i}$ and $\cE_{2}=\bigoplus L_{i}^{-1}$. Then
$\Ext^{2}(\cE_{1},\cE_{2})=\bigoplus_{i,j}H^{2}(L_{i}^{-1}\otimes L_{j}^{-1})$
and $\Ext^{2}(\cE_{2},\cE_{1})=\bigoplus_{i,j}H^{2}(L_{i}\otimes L_{j})$, and
the first Chern classes $\mp\bigl(c_{1}(L_{i})+c_{1}(L_{j})\bigr)$ occurring
here, for all $i$ and $j$, are differences of the kind in the hypothesis,
with $\varepsilon'=-\varepsilon$. A nondegenerate line bundle has cohomology only in
the degree equal to its index, so the index hypothesis gives the required
$\Ext^{2}$ vanishing. The sheaf $\cE_{1}$ has rank $m\ne0$, and the Chern
character hypothesis is that of \Cref{thm:rankzero}.
Since $e^{\lambda}$ and $e^{-\lambda}$ combine into $\cosh\lambda$ and
$\sinh\lambda$, which occupy disjoint degrees, the map splits by parity, and
the odd part is the annihilator of \Cref{lem:annihilator} tensored with the
sign vector. At $n=3$ the dimensions are read off from the proof of
\Cref{prop:explicitobject}, where the signs are absorbed in the same way:
the kernel consists of the vectors $(u_{i}^{-k}z)_{i}$ with $z$ in the mixed
block $V_{+}^{0,1}\otimes V_{-}^{0,1}$ for $k=3$, which change sign under
$u_{i}\mapsto-u_{i}$, and with $z$ in $\bigwedge^{2}V_{+}^{0,1}$ or
$\bigwedge^{2}V_{-}^{0,1}$ for $k=2$ or $k=4$, which do not.
\end{proof}

\begin{remark}[Uniformity over the configurations]
\label{rem:familyuniform}
By \Cref{cor:pteobstructed} the kernel has dimension at least $n^{2}$ for
every configuration of line bundles $L_{i}^{\pm1}$ whose virtual Chern
character lies on the Weil line and whose differences are nondegenerate of
index different from $2$. The deficiency is therefore not special to the six
line bundles of \Cref{prop:explicitobject}: \Cref{thm:rankzero} produces it
from the vanishing of the rank alone.

\Cref{thm:rankzero} excludes the economical choice described at the start of
this subsection. An object on which $H^{2}(\cO_{A})$ acts faithfully and which
realises a Weil class through semiregularity cannot have its Chern character
concentrated on the Weil line. It must carry a nonzero rank, or nonzero Chern
character in other degrees, so that $z\mapsto z\cdot\ch(\cE)$ is injective on
$H^{0,2}$; the degree-two component of that product is $\rk(\cE)\cdot z$, and
rank zero is the case in which it vanishes. Markman's sheaves have nonzero
rank, and this is essential to his method: the class he shows to be
deformation-invariant is the normalised character
$\exp(-c_{1}(\cE)/\rk\cE)\ch(\cE)$, which divides by the rank, and the objects
at $n=3$ are reflexive sheaves \cite{Mar25a,Mar25b}. This narrows the fourth
ingredient of \Cref{rem:deformation}: it is a single sheaf of positive rank
whose Chern character has a Weil component alongside other components, rather
than a difference of two sheaves whose classes cancel outside the Weil line.
\end{remark}

\Cref{thm:rankzero} is stated for a direct sum because there the element
$\xi_{z}$ is explicit. The mechanism applies to any perfect complex on which
$H^{2}(\cO_{A})$ acts faithfully, and in particular to the object used in
\Cref{sec:reduction}.

\begin{proposition}[Complexes with concentrated Chern character]
\label{prop:concentrated}
Let $A$ be of $(K,-1,n)$-Weil type and let $E$ be a nonzero perfect complex on
$A$ with $\ch_{k}(E)=0$ for $k\ne n$ and $\ch_{n}(E)=N\omega$ for some
$N\ne0$ and some nonzero $\omega\in\HW(A)$. If the natural map
\begin{equation}\label{eq:idmap}
  H^{2}(\cO_{A})\longrightarrow\Ext^{2}(E,E),\qquad z\longmapsto z\cdot\id_{E},
\end{equation}
is injective, then $\dim\ker\sigma_{E}\ge n^{2}$, so $E$ is not semiregular.
The map \eqref{eq:idmap} is injective whenever $E$ is represented by a virtual
sum of line bundles, in particular at every base point of
\Cref{cor:splitalgebraic}, \Cref{thm:quatweil} and \Cref{thm:everyn}.
\end{proposition}

\begin{proof}
For any perfect complex the semiregularity map satisfies
$\sigma_{E}(z\cdot\id_{E})=z\cdot\ch(E)$, since
$\operatorname{tr}\bigl(z\cdot\id_{E}\circ\exp\mathrm{At}(E)\bigr)
=z\cdot\operatorname{tr}\exp\mathrm{At}(E)=z\cdot\ch(E)$. Hence for $z$ in the
annihilator of $\omega$ in $H^{0,2}(A)$,
\[
  \sigma_{E}(z\cdot\id_{E})=N\,z\cdot\omega=0 ,
\]
and that annihilator has dimension $n^{2}$ by \Cref{lem:annihilator}. If
\eqref{eq:idmap} is injective, the classes $z\cdot\id_{E}$ for $z$ in a basis
of the annihilator are linearly independent, which gives the bound.

For the last sentence, write the class of $E$ in $K_{0}(A)$ as
$\sum_{i}\varepsilon_{i}[L_{i}]$ with $\varepsilon_{i}=\pm1$ and $L_{i}$ line
bundles, and realise it by the direct sum of the $L_{i}$ with
$\varepsilon_{i}=1$ and of the shifts $L_{i}[1]$ of those with
$\varepsilon_{i}=-1$. The map \eqref{eq:idmap} is then $z\mapsto(z,\dots,z)$
into $\bigoplus_{i}\Ext^{2}(L_{i},L_{i})=\bigoplus_{i}H^{2}(\cO_{A})$, which is
injective. At a split member $\omega_{1}$ is a polynomial in
$\beta,\widehat\beta,\ell$ by \Cref{thm:splitclosed}(iii), at a quaternionic
member it is a polynomial in divisor classes by \Cref{thm:quatweil}, and at a
product it is one by \Cref{thm:everyn}; in each case the Chern character is a
ring expression in classes of line bundles, so such a representation exists.
\end{proof}

\begin{proposition}[The class forces the size of the object]%
\label{prop:objectsize}
Let $A$ be of $(K,-1,n)$-Weil type, let $\omega$ be a nonzero element of the
Weil line and let $E$ be a perfect complex on $A$ with
\[
  \ch(E)=r+N\omega,\qquad r\in H^{0}(A,\QQ),\ N\ne0,
\]
and no other component. Then
\[
  \chi(E,E)\;=\;(-1)^{n}N^{2}\!\int_{A}\omega^{2},
\]
independently of $r$, and consequently
\[
  \sum_{i}\dim\Ext^{i}(E,E)\;\ge\;N^{2}\Bigl|\int_{A}\omega^{2}\Bigr|\;>\;0 .
\]
\end{proposition}

\begin{proof}
The Todd class of an abelian variety is $1$, so Riemann--Roch for the Ext
algebra reads $\chi(E,E)=\int_{A}\ch(E)^{\vee}\ch(E)$ with
$\ch_{k}^{\vee}=(-1)^{k}\ch_{k}$. Expanding,
\[
  \ch(E)^{\vee}\ch(E)
  = r^{2}+\bigl(1+(-1)^{n}\bigr)rN\omega+(-1)^{n}N^{2}\omega^{2},
\]
and of the three terms only the last has degree $4n$, so only it survives the
integral; the rank enters solely through $r^{2}$, which sits in degree zero.
For the inequality, $\chi(E,E)=\sum_{i}(-1)^{i}\dim\Ext^{i}(E,E)$, so the
total dimension is at least $|\chi(E,E)|$. Finally $\int_{A}\omega^{2}\ne0$.
Write $\omega=\alpha_{+}+\alpha_{-}$ with $\alpha_{\pm}\in\bigwedge^{2n}V_{\pm}$,
both nonzero by the proof of \Cref{lem:annihilator}. Each $\alpha_{\pm}$ is a
top form on its own eigenspace and squares to zero, so
$\omega^{2}=2\alpha_{+}\alpha_{-}$, which is a nonzero multiple of the
fundamental class.
\end{proof}

\begin{remark}[The constant]\label{rem:sizeconstant}
The bound concerns the whole Ext algebra and does not depend on the rank, so
increasing the rank does not weaken it. The constant $|\int\omega^{2}|$
depends on the normalisation of $\omega$; measured against $\eta$ it picks up
factors of $d$ and of the polarisation type, and no dimension-free growth
rate can be read off from it. The bound uses only that the Weil class is not
isotropic. Together with \Cref{prop:concentrated} it constrains an object
carrying the class in two ways. If the Chern character is concentrated on the
Weil line, the semiregularity map has a kernel of dimension at least $n^{2}$
unless $H^{2}(\cO_{A})$ acts on the object degenerately; and whatever the
rank, the Ext algebra has dimension at least $N^{2}|\int\omega^{2}|$.
\end{remark}

\begin{remark}[Consequences for the semiregularity reduction]
\label{rem:killsreduction}
\Cref{thm:semiregclosure} proves $\Sigma=\cD_{n,n}$, that is, the
algebraicity of the Weil class on the whole family, under the hypothesis that
$\sigma_{E}$ is injective for the complex $E$ of \Cref{prop:ktheory}, whose
Chern character is $N\omega_{1}$ in degree $n$ and zero elsewhere. This
complex satisfies the hypotheses of \Cref{prop:concentrated} at every base
point of \Cref{cor:splitalgebraic}, \Cref{thm:quatweil} and
\Cref{thm:everyn}, so there the hypothesis of \Cref{thm:semiregclosure} fails
for it, by $n^{2}$ dimensions.

This does not make \Cref{thm:semiregclosure} vacuous, for two reasons. First,
the hypothesis concerns a chosen perfect complex representing a class in
$K_{0}(A)_{\QQ}$, and distinct representatives of one class have distinct
$\Ext^{2}$ groups; what is ruled out is the representative that the
construction supplies. Second, \Cref{prop:concentrated} needs
\eqref{eq:idmap} to be injective, which is a real restriction for an object
of rank zero: it fails, for instance, for a sheaf supported at a point. A
representative escaping \Cref{prop:concentrated} must therefore
receive $H^{2}(\cO_{A})$ non-injectively, which for a rank-zero object means
that its support is degenerate. Of the two reductions of
\Cref{sec:reduction}, the degree bound of \Cref{thm:degreebound} is untouched
by anything here, and the semiregularity reduction of
\Cref{thm:semiregclosure} is unavailable for the complexes that the
constructions at these base points supply.
\end{remark}

\begin{figure}[t]
\centering
  \includegraphics[width=0.94\textwidth]{fig_annihilator}
\caption{The structure behind \Cref{lem:annihilator}, drawn at $n=3$.
$H^{0,2}(A)=\bigwedge^{2}H^{0,1}$ splits into three blocks by how many factors
come from $V_{+}^{0,1}$ and how many from $V_{-}^{0,1}$. The two outer blocks,
of dimension $\binom{n}{2}=3$ each, are carried injectively by multiplication
with the Weil class. The middle block, of dimension $n^{2}=9$ and spanned by
the mixed products, goes to zero, because each such product repeats a factor
of $\alpha_{+}$ and a factor of $\alpha_{-}$. Its dimension $n^{2}$ is the
deficit in \Cref{prop:concentrated}.}
\label{fig:annihilator}
\end{figure}


\subsection{The annihilator and the tangent space to the family}
\label{ssec:tangent}

The number $n^{2}$ of \Cref{lem:annihilator} is also the dimension of the
Weil family. We show in this subsection that the polarisation identifies the
annihilator with the tangent space to the family, canonically and in every
dimension. Two consequences follow. First, the deficiency of the
semiregularity map at an object of rank zero cannot be removed by a better
choice of object, since it is the family itself, seen inside $H^{0,2}$.
Second, the first-order Hodge locus of the Weil line is the Weil family, so
the directions that the semiregularity map cannot see are those along which
the class stays Hodge and along which it would have to be carried.

Write $V_{\pm}$ for the two eigenspaces of $K$ on $H^{1}(A,\CC)$, each of
dimension $2n$, and
\[
  V_{\pm}^{1,0}=V_{\pm}\cap H^{1,0},\qquad
  V_{\pm}^{0,1}=V_{\pm}\cap H^{0,1},
\]
each of dimension $n$ by \Cref{lem:balanced}. Let $E$ denote the
polarisation form on $H^{1}(A,\QQ)$, extended to $H^{1}(A,\CC)$. We use
throughout that $E$ vanishes on $V_{+}\times V_{+}$ and on $V_{-}\times V_{-}$
and pairs $V_{+}$ with $V_{-}$ perfectly, which is condition (iii) of
\Cref{def:weiltype} read on the eigenspaces, and that $E$ vanishes on
$H^{1,0}\times H^{1,0}$ and on $H^{0,1}\times H^{0,1}$, which is the first
Riemann bilinear relation. For the first statement, let $x,y\in V_{+}$. Then
$E(\sqrt{-d}\,x,y)=\sqrt{-d}\,E(x,y)$ from the eigenvalue, and
$E(\sqrt{-d}\,x,y)=-E(x,\sqrt{-d}\,y)=-\sqrt{-d}\,E(x,y)$ from
\Cref{def:weiltype}(iii), so $E(x,y)=0$. The same holds on $V_{-}$, and the
nondegeneracy of $E$ then forces the pairing between $V_{+}$ and $V_{-}$ to
be perfect. Together, the two vanishing statements say that of the four
blocks of $E$ on
$(V_{+}^{1,0}\oplus V_{+}^{0,1})\times(V_{-}^{1,0}\oplus V_{-}^{0,1})$ only
$V_{+}^{1,0}\times V_{-}^{0,1}$ and $V_{+}^{0,1}\times V_{-}^{1,0}$ are
nonzero, and each of these two is a perfect pairing of $n$-dimensional
spaces. This block structure is all that the proof of \Cref{thm:anntangent}
uses of $E$.

A first-order deformation of $A$ is a class in $H^{1}(A,T_{A})$, which for an
abelian variety is canonically $\Hom(H^{1,0},H^{0,1})$, since the tangent
bundle is trivial and $H^{1,0}$ is the dual of its space of sections. Such a
$v$ is extended to $H^{1}(A,\CC)$ by zero on $H^{0,1}$ and acts on
$H^{k}(A,\CC)=\bigwedge^{k}H^{1}(A,\CC)$ as the derivation it generates; we
write $v\lrcorner\,\xi$ for that action. A Hodge class $\xi$ of type $(p,p)$
remains of type $(p,p)$ to first order along $v$ exactly when
$v\lrcorner\,\xi=0$. We call $v$ \emph{polarised} when
$E(vx,y)+E(x,vy)=0$ for all $x,y\in H^{1,0}$, which is the condition that the
deformation preserve the polarisation. The polarised $v$ form a space of
dimension $n(2n+1)=\dim\cA_{2n}$.

\begin{proposition}[The first-order Hodge locus of the Weil line]
\label{prop:firstorder}
Let $A$ be of $(K,-1,n)$-Weil type and let $v$ be a polarised first-order
deformation of $A$. The following are equivalent.
\begin{enumerate}[label=\textup{(\roman*)},nosep]
\item $v\lrcorner\,\omega=0$ for every $\omega$ in the Weil line.
\item $v$ commutes with the action of $K$.
\item $v\bigl(V_{+}^{1,0}\bigr)\subseteq V_{+}^{0,1}$ and
  $v\bigl(V_{-}^{1,0}\bigr)\subseteq V_{-}^{0,1}$.
\end{enumerate}
The space of such $v$ has dimension $n^{2}$, and under
\Cref{prop:weildomain} it is the tangent space at $[A]$ to the Weil family
$\cD_{n,n}$. Each of the two rational generators of the Weil line already
cuts it out by itself.
\end{proposition}

\begin{proof}
The equivalence of (ii) and (iii) is immediate. Since $K$ acts on $V_{\pm}$
by the two scalars $\pm\sqrt{-d}$, commuting with $K$ means preserving $V_{+}$
and $V_{-}$. The extension of $v$ kills $H^{0,1}$ and takes values in
$H^{0,1}$, so preserving $V_{\pm}$ amounts to
$v(V_{\pm}^{1,0})\subseteq V_{\pm}\cap H^{0,1}=V_{\pm}^{0,1}$.

For (i) and (iii), let $\alpha_{+}$ generate $\bigwedge^{2n}V_{+}$ and write
$\alpha_{+}=e_{1}\wedge\cdots\wedge e_{2n}$ for a basis $e_{1},\dots,e_{2n}$
of $V_{+}$ adapted to $V_{+}=V_{+}^{1,0}\oplus V_{+}^{0,1}$. As a derivation,
\[
  v\lrcorner\,\alpha_{+}
  =\sum_{j=1}^{2n}e_{1}\wedge\cdots\wedge v(e_{j})\wedge\cdots\wedge e_{2n}.
\]
Split $v(e_{j})$ along $H^{1}(A,\CC)=V_{+}\oplus V_{-}$. The $V_{+}$
components contribute
$\operatorname{tr}\bigl(\pi_{+}\circ v|_{V_{+}}\bigr)\,\alpha_{+}$, and that
trace vanishes: the map $\pi_{+}\circ v|_{V_{+}}$ kills $V_{+}^{0,1}$ and
carries $V_{+}^{1,0}$ into $V_{+}^{0,1}$, so in the adapted basis its matrix
has zero diagonal blocks. Writing
$v(e_{j})_{-}=\sum_{k}c_{jk}f_{k}$ for a basis $f_{1},\dots,f_{2n}$ of
$V_{-}$, the $V_{-}$ components contribute
\[
  \sum_{j,k}\pm\,c_{jk}\;
  e_{1}\wedge\cdots\wedge\widehat{e_{j}}\wedge\cdots\wedge e_{2n}\wedge f_{k},
\]
and the $4n^{2}$ monomials appearing here are distinct members of the
standard basis of $\bigwedge^{2n}H^{1}(A,\CC)$, hence linearly independent.
So $v\lrcorner\,\alpha_{+}=0$ if and only if every $c_{jk}$ vanishes, which is
$v(V_{+})\subseteq V_{+}$. The same argument on $\alpha_{-}$ gives
$v\lrcorner\,\alpha_{-}=0$ if and only if $v(V_{-})\subseteq V_{-}$. The Weil
line spans $\bigwedge^{2n}V_{+}\oplus\bigwedge^{2n}V_{-}$ over $\CC$ by
\Cref{prop:weilline}, and each of the two rational generators of
\Cref{thm:explicitgens} has both components nonzero, so a single generator
already forces both conditions. This is (i) if and only if (iii).

For the dimension, let $v$ satisfy (iii), and write $v_{+}$ and $v_{-}$ for
its restrictions to $V_{+}^{1,0}$ and $V_{-}^{1,0}$. If $x$ and $y$ both lie
in $V_{+}^{1,0}$, then $v_{+}x$ and $y$ both lie in $V_{+}$, so both terms of
the polarisation condition vanish, and likewise if both lie in $V_{-}^{1,0}$.
For $x\in V_{+}^{1,0}$ and $y\in V_{-}^{1,0}$ the condition reads
$E(v_{+}x,y)=-E(x,v_{-}y)$. The left side is determined by $v_{+}$, and the
pairing between $V_{+}^{1,0}$ and $V_{-}^{0,1}$ is perfect, so for each $y$
there is exactly one $v_{-}y\in V_{-}^{0,1}$ satisfying this for all $x$, and
it depends linearly on $y$. Hence $v_{-}$ is determined by $v_{+}$, which may
be chosen freely, and the polarised deformations satisfying (iii) form a
space of dimension $\dim\Hom(V_{+}^{1,0},V_{+}^{0,1})=n^{2}$. By
\Cref{prop:weildomain} the family $\cD_{n,n}$ is an open subset of the
Grassmannian of $n$-planes in $V_{+}$, whose tangent space at
$V_{+}^{1,0}$ is $\Hom(V_{+}^{1,0},V_{+}/V_{+}^{1,0})=\Hom(V_{+}^{1,0},V_{+}^{0,1})$,
and the identification just made is that isomorphism.
\end{proof}

\begin{proposition}[Infinitesimal rigidity of divisor classes]%
\label{prop:rigidity}
Keep the notation above, put $P=V_{+}^{1,0}$ and $P'=V_{-}^{1,0}$ as in the
proof of \Cref{prop:weildomain}, so that $\bbar{P}=V_{-}^{0,1}$ and
$\bbar{P}'=V_{+}^{0,1}$, and write a class $\delta\in H^{1,1}$ in the four
blocks
\[
  a\ \text{on}\ P\otimes\bbar{P},\quad
  b\ \text{on}\ P'\otimes\bbar{P}',\quad
  c\ \text{on}\ P\otimes\bbar{P}',\quad
  e\ \text{on}\ P'\otimes\bbar{P},
\]
the first two carrying the norm character of $\Kx$ and the last two the
characters $\tau^{2}$ and $\bbar\tau^{\,2}$; reality makes $a$ and $b$
anti-hermitian and ties $e$ to $-\bbar{c}^{\,T}$. Identify the tangent space
$\Hom(P,\bbar{P}')$ to $\cD_{n,n}$ with the $n\times n$ matrices $V$ as in
\Cref{thm:anntangent}, and write $\operatorname{alt}$ for the alternating part
of a matrix, viewed in $\bigwedge^{2}$. Then contraction with a tangent
vector is
\begin{equation}\label{eq:phidelta}
  \Phi_{\delta}(V)
  \;=\;\bigl(\,Va+b^{T}V,\ \ \operatorname{alt}(Vc),\ \
       \operatorname{alt}(-V^{T}e)\,\bigr),
\end{equation}
with values in
$\bbar{P}'\otimes\bbar{P}\oplus\bigwedge^{2}\bbar{P}'\oplus\bigwedge^{2}\bbar{P}
=H^{0,2}$, and $\ker\Phi_{\delta}$ is the tangent space to the locus where
$\delta$ stays of type $(1,1)$. Consequently:
\begin{enumerate}[label=\textup{(\roman*)},nosep]
\item $\Phi_{\delta}=0$ if and only if $\delta$ is a real multiple of $\eta$.
  The rigid classes form a real line, which is the infinitesimal form of the
  statement that $\NS(A)_{\QQ}$ has rank one at the general member.
\item For the norm-character class with block $a$ and $b=c=e=0$,
  \[
    \dim\ker\Phi_{\delta}\;=\;n\,(n-\operatorname{rank}a),
  \]
  so such a class cuts a locus of codimension $n\operatorname{rank}a$, at
  least $n$.
\item A general class has $\ker\Phi_{\delta}=0$: to first order it
  rigidifies the family to a point.
\end{enumerate}
\end{proposition}

\begin{proof}
For \eqref{eq:phidelta}, a polarised deformation commuting with $K$ sends $P$
into $\bbar{P}'$ and $P'$ into $\bbar{P}$, and the polarisation condition
$E(vx,y)+E(x,vy)=0$ forces the second matrix to be $-V^{T}$, as in the proof
of \Cref{prop:firstorder}. Applying $v$ to the $(1,0)$ slot of each block and
collecting terms in $H^{0,2}$ gives the three displayed components; the sign
in the first comes from $\bbar{P}\wedge\bbar{P}'=-\bbar{P}'\wedge\bbar{P}$.

(i) If $\Phi_{\delta}=0$ then $\operatorname{alt}(Vc)=0$ for every $V$, which
for $n\ge2$ forces $c=0$, and likewise $e=0$. The first component gives
$Va=-b^{T}V$ for every $V$; taking $V$ to run over the matrix units shows
$a$ is a scalar $\lambda$ and $b^{T}=-\lambda$. Anti-hermitian scalars are
the purely imaginary ones, so the solutions form a real line. It contains
$\eta$, whose blocks are $a=i$, $b=-i$, and the two therefore agree.

(ii) With $b=c=e=0$ the condition is $Va=0$, that is, the rows of $V$
annihilate the column space of $a$; the solution space has dimension
$n(n-\operatorname{rank}a)$.

(iii) By (ii), the class with $a$ of rank $n$ (for instance $a=i$) and
$b=c=e=0$, which is real since $a$ is anti-hermitian, has
$\ker\Phi_{\delta}=0$. Since $\Phi_{\delta}$ depends linearly on $\delta$, the
classes with $\ker\Phi_{\delta}\ne0$ are cut out by the vanishing of the
maximal minors of $\Phi_{\delta}$ and form a proper real algebraic subset of
the real classes of type $(1,1)$, which gives (iii).
\end{proof}

\begin{remark}[The divisorial supply]\label{rem:rigiditywhy}
\Cref{prop:rigidity}(ii) explains why the divisorial supply of
\Cref{thm:divcriterion} cannot be enlarged by a better choice of class. Any
norm-character class other than a multiple of $\eta$ has
$\operatorname{rank}a\ge1$ and so cuts a locus of codimension at least $n$;
that locus is proper, and by (i) only the real multiples of $\eta$ cut
nothing. Each further class therefore yields a subvariety and never an open
set, which is \Cref{prop:divfails} with the codimension made explicit.
\end{remark}

\begin{figure}[htbp]
\centering
  \includegraphics[width=0.46\textwidth]{fig_fibres}
\caption{The geometry behind \Cref{lem:annihilator}, drawn schematically for
the real shadow of the hermitian form. The ruled surface is the null cone of
$H$; the two transverse planes are the conjugate isotropic subspaces $V_{+}$
and $V_{-}$, which conjugation interchanges, so that neither is defined over
$\QQ$ and only their sum is. The segments join a vector of one to its
conjugate in the other, and the strip they sweep stands for the mixed block
$V_{+}^{0,1}\otimes V_{-}^{0,1}$, of dimension $n^{2}$, which multiplication
by the Weil class sends to zero and which is the deficit in
\Cref{prop:concentrated}.}
\label{fig:fibres}
\end{figure}

\begin{theorem}[The annihilator is the tangent space]\label{thm:anntangent}
Let $A$ be of $(K,-1,n)$-Weil type and let $\omega$ be a rational generator of
the Weil line. The polarisation induces an isomorphism
\[
  T_{[A]}\cD_{n,n}\;\xrightarrow{\ \sim\ }\;
  \operatorname{Ann}_{H^{0,2}(A)}(\omega),
\]
canonical in the sense that it depends on no choice of basis and on no choice
of $\omega$ on the line. Explicitly, if $p_{1},\dots,p_{n}$ is a basis of
$V_{+}^{1,0}$ and $q_{1},\dots,q_{n}$ the basis of $V_{-}^{0,1}$ dual to it
under $E$, then $\varphi\in\Hom(V_{+}^{1,0},V_{+}^{0,1})$ corresponds to
\[
  \sum_{a=1}^{n}q_{a}\wedge\varphi(p_{a})\;\in\;
  V_{-}^{0,1}\otimes V_{+}^{0,1}\;\subset\;\textstyle\bigwedge^{2}H^{0,1}
  =H^{0,2}(A).
\]
\end{theorem}

\begin{proof}
The form $E$ pairs $V_{+}$ with $V_{-}$ perfectly and vanishes on
$V_{+}^{1,0}\times V_{-}^{1,0}$, the latter because both spaces lie in
$H^{1,0}$. Hence $E$ restricted to $V_{+}^{1,0}\times V_{-}$ kills
$V_{-}^{1,0}$ and descends to a pairing
\[
  V_{+}^{1,0}\times V_{-}^{0,1}\longrightarrow\CC .
\]
This pairing is perfect: both spaces have dimension $n$, and a vector of
$V_{+}^{1,0}$ orthogonal to $V_{-}^{0,1}$ is orthogonal to all of $V_{-}$,
hence zero. It gives a canonical isomorphism
$\bigl(V_{+}^{1,0}\bigr)^{*}\cong V_{-}^{0,1}$, and therefore
\[
  T_{[A]}\cD_{n,n}
  =\Hom\bigl(V_{+}^{1,0},V_{+}^{0,1}\bigr)
  =\bigl(V_{+}^{1,0}\bigr)^{*}\otimes V_{+}^{0,1}
  \;\cong\;V_{-}^{0,1}\otimes V_{+}^{0,1},
\]
the first equality being the tangent space computation at the end of the proof
of \Cref{prop:firstorder}. Since $H^{0,1}=V_{+}^{0,1}\oplus V_{-}^{0,1}$, the
map $q\otimes u\mapsto q\wedge u$ identifies $V_{-}^{0,1}\otimes V_{+}^{0,1}$
with the middle summand of
$\bigwedge^{2}H^{0,1}$, and that summand is the annihilator by
\Cref{lem:annihilator}. The displayed formula is the composite of these
isomorphisms written in a dual pair of bases, so it is independent of the
basis chosen; and \Cref{lem:annihilator} gives the same subspace for every
$\omega$ on the line, so the isomorphism does not depend on $\omega$ either.
\end{proof}

\begin{figure}[tbp]
\centering
  \includegraphics[width=\textwidth]{fig_tangent}
\caption{\Cref{thm:anntangent}. The polarisation has only two nonzero blocks
on $V_{+}\times V_{-}$: it pairs $V_{+}^{1,0}$ with $V_{-}^{0,1}$ and
$V_{+}^{0,1}$ with $V_{-}^{1,0}$, each perfectly, and the other two blocks
vanish by the first Riemann relation. The upper right block turns the tangent
space $\Hom(V_{+}^{1,0},V_{+}^{0,1})$ to the Weil family into
$V_{-}^{0,1}\otimes V_{+}^{0,1}$, which is the middle block of
$\bigwedge^{2}H^{0,1}$ drawn in \Cref{fig:annihilator} and the annihilator of
the Weil line. The two $n^{2}$-dimensional spaces are therefore the same
subspace of $H^{0,2}$, and the directions that the semiregularity map cannot
see are the directions of the family along which the class stays Hodge.}
\label{fig:tangent}
\end{figure}


\begin{corollary}[The kernel contains the tangent space to the family]
\label{cor:kernelfamily}
Let $A$ be of $(K,-1,n)$-Weil type and let $E$ be a nonzero perfect complex on
$A$ satisfying the hypotheses of \Cref{prop:concentrated}, that is, with Chern
character concentrated on the Weil line and with \eqref{eq:idmap} injective.
Then $\ker\sigma_{E}$ contains a subspace canonically isomorphic to
$T_{[A]}\cD_{n,n}$, namely the image of the composite
\[
  T_{[A]}\cD_{n,n}\;\cong\;\operatorname{Ann}_{H^{0,2}}(\omega)
  \;\hookrightarrow\;H^{2}(\cO_{A})\;\xrightarrow{\;z\mapsto z\cdot\id_{E}\;}\;
  \Ext^{2}(E,E).
\]
By \Cref{prop:firstorder} that subspace is also the tangent space to the
first-order Hodge locus of the class $\ch_{n}(E)$.
\end{corollary}

\begin{proof}
Immediate from \Cref{thm:anntangent}, \Cref{prop:concentrated} and
\Cref{prop:firstorder}.
\end{proof}

\begin{theorem}[The weakened criterion is unavailable for concentrated
Chern characters]\label{thm:weakfails}
Let $A$ be of $(K,-1,n)$-Weil type and let $E$ be a nonzero perfect complex on
$A$ satisfying the hypotheses of \Cref{prop:concentrated}. Let
\[
  \operatorname{ev}_{E}\colon HH^{2}(A)\longrightarrow\Ext^{2}(E,E)
\]
be the evaluation on $E$ of the natural transformations $\id\to\id[2]$ of
$D^{b}(A)$, as in \cite[Section 11.4]{Mar25b}. Then
\[
  \operatorname{ev}_{E}\bigl(H^{2}(\cO_{A})\bigr)\cap\ker\sigma_{E}
  =\Delta:=\bigl\{\,z\cdot\id_{E}\;:\;z\in\Ann_{H^{0,2}}(\omega)\,\bigr\},
\]
a subspace of dimension $n^{2}$. In particular $\sigma_{E}$ is not injective
on the image of $\operatorname{ev}_{E}$, so $E$ satisfies neither the
hypothesis of \cite[Lemma 11.3]{Mar25b} nor the weakened hypothesis of
\cite[Question 11.4]{Mar25b}.
\end{theorem}

\begin{proof}
The Hochschild cohomology $HH^{*}(A)=\Ext^{*}_{A\times A}(\cO_{\Delta},\cO_{\Delta})$
is the algebra of natural transformations of the identity functor, and the
Hochschild--Kostant--Rosenberg isomorphism \cite{Cal05,BF08} writes
$HH^{2}(A)\cong H^{2}(\cO_{A})\oplus H^{1}(A,T_{A})\oplus H^{0}(A,\bigwedge^{2}T_{A})$,
with no correction from the Todd class because $T_{A}$ is trivial. The first
summand is the image of $H^{2}(\cO_{A})=\Ext^{2}(\cO_{A},\cO_{A})$ under
pullback to $A\times A$, and the natural transformation it defines sends a
complex $F$ to the composite $F=F\otimes\cO_{A}\xrightarrow{\;\id\otimes z\;}
F\otimes\cO_{A}[2]=F[2]$, that is, to $z\cdot\id_{F}$. This is the statement
that the action of $HH^{*}(A)$ on $\Ext^{*}(F,F)$ restricts on $H^{*}(\cO_{A})$
to the $\cO_{A}$-module structure. Concretely, a class $z\in H^{2}(\cO_{A})$
is represented by a $2$-extension of $\cO_{A}$ by itself; its pullback
$\pi_{1}^{*}z$ to $A\times A$ is a $2$-extension of $\cO_{A\times A}$ by
itself, and tensoring with $\cO_{\Delta}$ gives a $2$-extension of
$\cO_{\Delta}$ by itself, which is the element of $HH^{2}(A)$ in question. The
Fourier--Mukai transform with kernel $\cO_{\Delta}$ is the identity, and the
transform with kernel this extension is the extension of the identity by
itself twisted by $z$, so its value on $F$ is $z\cdot\id_{F}$. Markman uses
this identification when he writes the left vertical arrow of the diagram in
the proof of \cite[Lemma~11.3]{Mar25b} as the HKR isomorphism followed by
contraction with $\ch(E)$. On the summand $H^{2}(\cO_{A})$ contraction is the
cup product $z\mapsto z\cdot\ch(E)$, which is $\sigma_{E}(z\cdot\id_{E})$ by
the computation in the proof of \Cref{thm:rankzero}.

Hence $\operatorname{ev}_{E}(z)=z\cdot\id_{E}$ for $z\in H^{2}(\cO_{A})$, and
$\sigma_{E}(z\cdot\id_{E})=z\cdot\ch(E)=N\,z\cdot\omega$ by the proof of
\Cref{prop:concentrated}, which vanishes exactly for
$z\in\Ann_{H^{0,2}}(\omega)$. So
$\operatorname{ev}_{E}(H^{2}(\cO_{A}))\cap\ker\sigma_{E}=\Delta$, and
$\dim\Delta=\dim\Ann_{H^{0,2}}(\omega)=n^{2}$ by \Cref{lem:annihilator} and
the injectivity of \eqref{eq:idmap}. Since $\Delta\ne0$, the restriction of
$\sigma_{E}$ to the image of $\operatorname{ev}_{E}$ has a nonzero kernel.
The hypothesis of \cite[Lemma 11.3]{Mar25b} is that
$\ker\operatorname{ev}_{E}$ equals the kernel of contraction with $\ch(E)$ and
that $\operatorname{ev}_{E}$ is surjective; the two conditions together would
make $\sigma_{E}$ injective on all of $\Ext^{2}(E,E)$, contradicting
\Cref{prop:concentrated}. The weakened hypothesis of
\cite[Question 11.4]{Mar25b} is injectivity of $\sigma_{E}$ on the image,
which we have just shown to fail.
\end{proof}

\begin{theorem}[A line bundle summand off the polarisation line obstructs]
\label{thm:linebundle}
Let $A$ be of $(K,-1,n)$-Weil type with $n\ge2$ and let $T=T_{[A]}\cD_{n,n}$.
Let $E$ be a perfect complex on $A$ with a direct summand $L[p]$, $L$ a line
bundle, and suppose that $\ch(E)$ remains of Hodge type to first order along
$\cD_{n,n}$, that is $v\lrcorner\,\ch(E)=0$ for every $v\in T$. If $c_{1}(L)$
is not a real multiple of $\eta$, then there is $v\in T$ with
\[
  \operatorname{ev}_{E}(v)\ne0
  \qquad\text{and}\qquad
  \sigma_{E}\bigl(\operatorname{ev}_{E}(v)\bigr)=0 .
\]
Consequently $E$ is not semiregular, does not satisfy the weakened hypothesis
of \cite[Question 11.4]{Mar25b}, and is obstructed to first order along $v$.
The same holds for a twisted complex, with $c_{1}(L)$ replaced by the twisted
first Chern class of the summand, which is $c_{1}(L)$ shifted by the class
common to all summands.
\end{theorem}

\begin{proof}
The class $c_{1}(L)$ is a real class of type $(1,1)$ and not a real multiple
of $\eta$, so by \Cref{prop:rigidity}(i) the contraction
$\Phi_{c_{1}(L)}\colon T\to H^{0,2}(A)$ is not zero: there is $v\in T$ with
$v\lrcorner\,c_{1}(L)\ne0$. On the summand $H^{1}(A,T_{A})$ of $HH^{2}(A)$ the
evaluation map is contraction with the Atiyah class,
$\operatorname{ev}_{E}(v)=v\lrcorner\,\mathrm{At}(E)$. This is the top row of
the diagram of Buchweitz and Flenner reproduced as \cite[(2.1)]{Mar25b}; its
identification with the Hochschild evaluation is in \cite{BF08}, and Markman
uses it in the proof of \cite[Lemma~11.3]{Mar25b}. The Atiyah class of a
direct sum is the direct sum of the Atiyah classes, and that of $L[p]$ is
$c_{1}(L)$, so the
component of $\operatorname{ev}_{E}(v)$ in the summand
$\Ext^{2}(L[p],L[p])=H^{0,2}(A)$ of $\Ext^{2}(E,E)$ is
$v\lrcorner\,c_{1}(L)\ne0$. On the other hand the same diagram gives
$\sigma_{E}(\operatorname{ev}_{E}(v))=v\lrcorner\,\ch(E)$, which vanishes by
hypothesis. So $\operatorname{ev}_{E}(v)$ is a nonzero element of
$\ker\sigma_{E}\cap\image\operatorname{ev}_{E}$; this is the failure of both
hypotheses, and $v\lrcorner\,\mathrm{At}(E)\ne0$ is the statement that $E$
does not lift to a first-order deformation of the pair along $v$, as recalled
in \cite[Section 2]{Mar25b}. Twisting by a class $B$ replaces the Atiyah
class of every rank one summand by its shift by $B$ and $\ch(E)$ by the
twisted Chern character, and nothing else in the argument changes.
\end{proof}

\begin{remark}[Scope of the kernel computation]
\label{rem:whatmeasures}
The semiregularity criterion converts a statement about the Chern character
into a statement about the object by requiring $\sigma_{E}$ to be injective.
\Cref{cor:kernelfamily} says that for an object whose Chern character is
concentrated on the Weil line this requirement fails, and that the failure is
measured by the tangent space to the family across which the class is to be
transported, in every dimension and with no hypothesis beyond
\eqref{eq:idmap}. This concerns the size and the position of the kernel. By
itself it does not say that the obstruction classes are nonzero, nor that
the Weil class fails to deform. For the explicit object of
\Cref{prop:explicitobject} the obstruction classes are computed in
\Cref{prop:splitobstruction}: they are nonzero in every direction of the
family normal to the split locus, they lie in $\ker\sigma_{E}$, and they meet
the subspace of \Cref{cor:kernelfamily} only in zero. The Weil class deforms
along the family, by \Cref{prop:firstorder}, while that object does not.

The same computation bears on a weaker criterion, which it settles only in
part. Markman asks in \cite[Question 11.4]{Mar25b} whether the
semiregularity theorem survives when $\sigma_{E}$ is required to be injective
only on the image of the evaluation map $\operatorname{ev}_{E}$ rather than
on all of $\Ext^{2}(E,E)$, and records in his Section 12 that an affirmative
answer would carry the method to higher dimensions and to fields of larger
degree. By \Cref{thm:weakfails} the classes of \Cref{cor:kernelfamily} lie in
that image, so for concentrated Chern characters the weakened hypothesis
fails as well, in every dimension and by the same $n^{2}$ dimensions. By
\Cref{thm:linebundle} it fails for every object with a line bundle summand
off the polarisation line, whatever the rank and at every member of the
family, and by \Cref{cor:noline} this excludes every sum of line bundles.
For the explicit split object of \Cref{prop:explicitobject} the failure is
complete: \Cref{prop:evaluation} shows that every one of the fifteen
dimensions of $\ker\sigma_{E}$ is the value of $\operatorname{ev}_{E}$ on a
Hochschild class that already preserves $\ch(E)$ to first order.

The next question is whether an object of positive rank can satisfy the
weakened hypothesis without satisfying the original one. Markman's own
sheaves are of positive rank and are not of the concentrated shape
(\Cref{rem:familyuniform}); for them he avoids the question by descending to
a quotient on which $\sigma_{E}$ is injective on an invariant subspace
containing the image of $\operatorname{ev}_{E}$
\cite[Proposition~11.5]{Mar25b}. In his Section 12 he proposes a rank one
complex on a principally polarised abelian fourfold, built from $(d+9)/2$
translates of the theta divisor, as a candidate for the weakened criterion in
dimension eight. \Cref{thm:candidatefails}
shows that the candidate fails, for every discriminant, at every isolated
double point of its support.

Voisin's account of the subject \cite{Voi26} makes the same point
qualitatively, for the variational Hodge conjecture in general: semiregular
objects are hard to construct and too restricted to carry a general argument.
\Cref{cor:kernelfamily} gives this observation a quantitative form for one
family of classes, in which the restriction amounts to $n^{2}$ dimensions and
those dimensions are the tangent directions of the family. Totaro's account
\cite{Tot26} describes Markman's method as the construction of enough
semiregular sheaves to prove the conjecture for abelian fourfolds and
fivefolds, and records that it needed the semiregularity theorem for twisted
sheaves, which Pridham supplied \cite{Pri24}.
\end{remark}

